\documentclass[11pt]{article}

\usepackage[margin=1in]{geometry}
\usepackage{amsmath,amssymb,amsthm,mathtools}
\usepackage{enumitem}
\usepackage{booktabs}
\usepackage{tabularx}
\usepackage{microtype}
\usepackage[hidelinks]{hyperref}
\usepackage{xurl}

\newtheorem{definition}{Definition}
\newtheorem{lemma}{Lemma}
\newtheorem{remark}{Remark}

\title{Challenge Surface Hooks for Successor Maintenance:\\[0.5ex]
\large Piece-to-Segment Realization and Shared Surface Bindings}
\author{}
\date{Unified archive note v0.04 (2026-03-21; archive v487)}

\begin{document}
\maketitle

\begin{abstract}
A downstream normal form already exposes auditable sentence segments.
A challenge coverage certificate already says which semantic pieces count as covered for which audiences and which hooks remain support-only.
One practical ambiguity remains anyway: when a terse emitted sentence carries several semantic pieces, where exactly is each piece realized on the surface, and when are those spans shared rather than piece-exclusive?
This note isolates that surface-realization layer.
It defines a compact \emph{challenge surface-hook ledger} mapping each covered sentence-family piece to exact emitted segments, carried claim classes, and a small binding-kind token recording whether those segments are exclusive, shared, or mixed.
Exact narrow terminal ownership is imported from a companion terminal-witness ledger.
\end{abstract}

\paragraph{Non-redundancy note.}
Synthesis~31 owns the downstream normal form and sentence segmentation, Synthesis~37 owns question-keyed semantic pieces, Synthesis~39 owns sentence-family coverage across audiences, Synthesis~41 owns exact piece-hook terminal ownership, Synthesis~54 owns the answer-side terminal citation normal form, Synthesis~74 owns the paired terminal citation discipline, Synthesis~75 owns the paired cutover rule, and Synthesis~17 owns the concrete worked instantiation.
This note owns only the exact piece-to-segment realization layer once coverage and the later-paper stop posture are already fixed. Later notes should reopen it only when the exact piece-to-span binding or shared-vs-exclusive surface claim is itself live.
\paragraph{Scope.}
This is not a new route, stop profile, branch map, coverage certificate, or terminal-witness ledger.
It is a narrow publication-interface note about saying \emph{where} a covered piece appears on the emitted surface.

\paragraph{Imports.}
The downstream normal form is imported from Synthesis~31.\cite{series:downstreamnormalforms}
Question-keyed semantic pieces are imported from Synthesis~37.\cite{series:challengestops}
Sentence-family coverage certificates are imported from Synthesis~39.\cite{series:challengecoverage}
Piece-hook terminal ownership is imported from Synthesis~41.\cite{series:challengeterminalwitnesses}
The answer-side terminal citation normal form and its paired reuse/cutover discipline are imported from Synthesis~54 and Synthesis~74--75 as the default later-paper stop rule around this local surface-binding owner.\cite{series:challengeanswerreviewmenus,series:pairedterminalcitationdiscipline,series:pairedterminalcutoverrules}
The concrete worked instantiation is imported from Synthesis~17.\cite{series:workedexample}
\paragraph{Exports.}
This note exports (i) challenge surface-hook ledgers, (ii) a hook-lookup rule tying coverage entries to exact piece realizations, (iii) a shared-vs-exclusive binding-kind rule, and (iv) a local-reopen rule saying that later notes may cite one maintained surface-hook ledger only when the exact piece-to-span binding or shared-vs-exclusive surface claim is itself live.
Otherwise later terse papers should default to the answer-side terminal citation normal form in Synthesis~54 together with the paired terminal discipline and paired cutover rule in Synthesis~74--75. Exact narrow terminal ownership remains outside scope here and belongs to Synthesis~41.\cite{series:challengeterminalwitnesses,series:challengeanswerreviewmenus,series:pairedterminalcitationdiscipline,series:pairedterminalcutoverrules}
\section{Why surface hooks deserve their own note}
Sentence segmentation alone says where the sentence breaks.
Coverage alone says which piece counts as covered for which audience.
A companion terminal-witness ledger says which narrow terminal owner belongs to that piece hook.
But neither object, by itself, cleanly answers the editorial question a referee or maintainer will still ask: which emitted span realizes \emph{this} piece, and is that span shared with another piece or not?
If that answer is left implicit, terse worked examples become decorative exactly where they should be most auditable: in the many-to-many interface between semantic pieces and public wording.
So exact surface realization deserves one small ledger of its own.

\begin{definition}[Challenge surface-hook ledger]
Fix one concrete base$\to$successor comparison.
A \emph{challenge surface-hook ledger} is a finite family
\[
\Sigma=\{(h,f,p,b,S,C)\}
\]
whose entries each contain a hook key $h$, a sentence family $f$, a semantic piece key $p$, a binding-kind token $b\in\{\textsf{exclusive},\textsf{shared},\textsf{mixed}\}$, a finite nonempty set $S$ of emitted segment keys, and a finite nonempty set $C$ of carried claim classes.
\end{definition}

\begin{definition}[Hook-lookup rule]
A surface-hook ledger satisfies the \emph{hook-lookup rule} if every hook key named by a companion coverage certificate resolves to one unique ledger entry with the same sentence family and piece key.
\end{definition}

\begin{definition}[Binding-kind rule]
A surface-hook entry satisfies the \emph{binding-kind rule} if
\begin{enumerate}[leftmargin=*]
\item $b=\textsf{exclusive}$ means all listed segments are used only by this piece hook in the worked sentence family,
\item $b=\textsf{shared}$ means all listed segments are also reused by at least one other piece hook in the same sentence family, and
\item $b=\textsf{mixed}$ means the listed set contains both shared and piece-exclusive segments.
\end{enumerate}
\end{definition}

\begin{definition}[Claim-carry rule]
A surface-hook entry satisfies the \emph{claim-carry rule} if every carried claim class in $C$ is carried by at least one listed emitted segment in $S$.
Exact terminal discharge for those classes is then checked in the companion terminal-witness ledger rather than here.
\end{definition}

\begin{lemma}[Surface hooks prevent hidden span sharing]
Assume a challenge surface-hook ledger satisfies the hook-lookup rule, the binding-kind rule, and the claim-carry rule.
Then a later note may reuse the sentence family without silently changing where one covered semantic piece is realized on the surface or whether that realization is shared with another piece.
\end{lemma}

\begin{proof}
The hook-lookup rule prevents audience-coverage entries from hand-waving at unspecified spans.
The binding-kind rule prevents shared realizations from being mistaken for piece-exclusive wording and vice versa.
The claim-carry rule prevents a span from being named as a realization of a piece it does not even carry on the sentence surface.
So the ledger records exact surface realization without making either the coverage note or the terminal-witness note own segment-level details.
\end{proof}

\begin{table}[t]
\centering
\small
\begin{tabularx}{\linewidth}{@{}l l X@{}}
\toprule
Hook kind & Typical use & Worked meaning \\
\midrule
Exclusive & one piece, one span family & public-status continuity half; public-status closure half \\
Shared & all listed spans reused elsewhere & changed-family identity carried by the same fallback spans as numeric deltas \\
Mixed & some shared, some exclusive & compact diff numeric deltas use two shared fallback spans plus one total-only span \\
\bottomrule
\end{tabularx}
\caption{Surface hooks say exactly where one covered piece appears on the emitted surface, including whether those spans are shared.}
\label{tab:surfacehooks}
\end{table}

\section{Worked example pointer}
The maintained worked bundle now ships \path{example_successor_challenge_surface_hooks.json}.\cite{series:workedexample}
It records four hooks for the canned Case-B successor.
The public-status sentence has two exclusive hooks: one for the continuity half and one for the closure half.
The compact diff sentence has one mixed hook for numeric deltas and one shared hook for changed-family identity, making explicit that the two fallback spans simultaneously realize both pieces while the total-summary span realizes only the numeric piece.
Exact narrow terminal ownership for those hooks is left to the companion terminal-witness ledger.
That keeps the normal-form note about sentence segmentation, the coverage note about audience obligations, the terminal-witness note about narrow stop ownership, and this note about exact piece-to-segment realization sharply separated.\cite{series:downstreamnormalforms,series:challengecoverage,series:challengeterminalwitnesses}

\paragraph{Why this helps the archive.}
The series can now state, without bloating the public prose, not only \emph{that} a terse sentence covers a piece, but also exactly \emph{where} that piece appears on the surface and whether that wording is shared with another piece.
That makes the worked example more referee-grade without making it longer.

\paragraph{A minimal public challenge-surface-hook card.}
\begin{table}[t]
\centering
\small
\begin{tabularx}{\linewidth}{@{}lX@{}}
\toprule
Field & Minimal public answer \\
\midrule
Public-status continuity hook & \nolinkurl{public_status_sentence_continuity_piece_surface} binds \nolinkurl{continuity_piece} in \nolinkurl{public_status_sentence} to the exclusive segment \nolinkurl{status_continuity_half} carrying the \nolinkurl{continuity_decision} claim class \\
Compact-diff delta hook & \nolinkurl{compact_numeric_diff_diff_delta_piece_surface} binds \nolinkurl{diff_delta_piece} in \nolinkurl{compact_numeric_diff} to the mixed segment set \nolinkurl{fallback_obs_delta}, \nolinkurl{fallback_summary_delta}, and \nolinkurl{total_summary_delta}, carrying the \nolinkurl{numeric_delta} claim class \\
Compact-diff changed-family hook & \nolinkurl{compact_numeric_diff_changed_family_piece_surface} binds \nolinkurl{changed_family_piece} in \nolinkurl{compact_numeric_diff} to the shared segment pair \nolinkurl{fallback_obs_delta} and \nolinkurl{fallback_summary_delta}, carrying the \nolinkurl{changed_family_identity} claim class \\
Selection rule & Stop here only while the live issue is where one already-covered piece is realized on the emitted surface and whether that realization is exclusive, shared, or mixed; do not widen to exact narrow terminal owners or proving bundles unless those objects have become the honest moving part \\
Left/right boundary & Reopen Synthesis~39 when the live issue is whether the sentence family covers the piece at all; widen to Synthesis~41 when the exact narrow terminal owner is live, and to Synthesis~42--44 when the honest issue is the proving bundle rather than the surface binding \\
Paired-terminal boundary & Once the exact piece-to-span binding is fixed, later terse papers should usually stop at Synthesis~54 together with Synthesis~74--75 rather than reopen this note merely to restate a stable answer-side handoff or paired cutover judgment \\
\bottomrule
\end{tabularx}
\caption{A minimal public challenge-surface-hook card. This is the smallest surface-binding object later terse papers can quote directly when the live issue is where one already-covered piece appears on the emitted sentence and whether that wording is exclusive, shared, or mixed.}
\label{tab:minimal-public-challenge-surface-hook-card}
\end{table}

This card is intentionally non-decorative: it pins the actual worked hook, segment, and claim-class keys together with the stop rule for one maintained piece-to-span binding, without silently re-owning the earlier family-level coverage claim or the later exact terminal-owner and proving-bundle objects.\cite{series:workedexample}

\paragraph{A forced public challenge-surface-hook matrix.}
\begin{table}[t]
\centering
\small
\begin{tabularx}{\linewidth}{@{}lX@{}}
\toprule
Field & Forced public answer \\
\midrule
Public-status continuity floor & \nolinkurl{public_status_sentence_continuity_piece_surface} still forces the same surface binding only while the same sentence family and piece key agree, the binding kind stays \textsf{exclusive}, the same segment \nolinkurl{status_continuity_half} is used, and that segment still carries \nolinkurl{continuity_decision} \\
Compact-diff delta floor & \nolinkurl{compact_numeric_diff_diff_delta_piece_surface} still forces the same surface binding only while the same sentence family and piece key agree, the binding kind stays \textsf{mixed}, the same three segments remain bound, and those segments still carry the \nolinkurl{numeric_delta} claim class \\
Changed-family floor & \nolinkurl{compact_numeric_diff_changed_family_piece_surface} still forces the same surface binding only while the same sentence family and piece key agree, the binding kind stays \textsf{shared}, the same fallback segment pair remains bound, and those segments still carry \nolinkurl{changed_family_identity} \\
Coverage boundary floor & The shortcut remains honest only while the companion coverage entry still says the piece is \textsf{surface}-realized for this sentence family; if the maintained piece becomes support-only or ceases to be covered here at all, the surface-hook citation has already gone stale \\
Terminal/proving boundary floor & This note still owns the same stop only while exact narrow terminal ownership and proving bundles remain imported neighbor objects; once the live issue becomes the exact terminal owner or proving burden, the honest stop has moved to Synthesis~41--44 \\
Staleness trigger floor & This matrix goes stale as soon as the sentence family changes, the piece key changes, the binding kind flips, a listed segment is added, removed, or rebound, a carried claim class drifts, or the maintained hook stops being the same exact piece-to-span realization \\
Escalation pointer & Stop here only while one maintained hook still forces the same emitted spans; reopen Synthesis~39 when the live issue is the coverage claim itself, widen to Synthesis~41 when the exact narrow owner is live, and widen to Synthesis~42--44 when the dispute has become the proving bundle rather than the surface binding \\
\bottomrule
\end{tabularx}
\caption{A forced public challenge-surface-hook matrix. This is the smallest surface-binding object later terse papers can quote directly when the live issue is whether one maintained hook still forces the same exact piece-to-span realization rather than merely which hook exists in the abstract.}
\label{tab:forced-public-challenge-surface-hook-matrix}
\end{table}

This matrix is intentionally non-decorative: it pins the actual worked hook, segment, and claim-class keys that still force one exact piece-to-span realization, together with the conditions under which that forcing remains honest or goes stale, without silently re-owning the earlier coverage claim or the later exact terminal-owner and proving-bundle objects.\cite{series:workedexample}

\paragraph{One tiny challenge-surface-hook drill.}
For the worked compact diff sentence, \nolinkurl{compact_numeric_diff_changed_family_piece_surface} still forces the same surface binding only while the same changed-family piece is carried by the same shared fallback pair \nolinkurl{fallback_obs_delta} and \nolinkurl{fallback_summary_delta}. If the live question widens from ``does this maintained hook still bind the same piece to the same spans?'' to ``which exact narrow terminal owner discharges that changed-family claim?'' the honest stop is no longer this note but Synthesis~41. If the live question stays about surface realization but the piece becomes support-only, the shared pair changes, or the carried claim class is rebound, the matrix has gone stale and a later terse paper must reopen the exact surface-hook owner rather than keep quoting the old piece-to-span binding. Only once the exact surface binding is fixed and the remaining issue becomes the later answer-side handoff or paired terminal cutover should the paper leave this note for Synthesis~54 or Synthesis~74--75.\cite{series:workedexample}

\begin{thebibliography}{99}\small

\bibitem{series:downstreamnormalforms}
Anonymous.
\newblock Downstream Normal Forms for Successor Maintenance: Summary, Support, Citation, Quote, and Clause Discipline.
\newblock Series synthesis note (Synthesis~31), 2026. (This archive.)

\bibitem{series:challengestops}
Anonymous.
\newblock Challenge Stop Profiles for Successor Maintenance: Question Families, Semantic Pieces, and Approved Terminal Fields.
\newblock Series synthesis note (Synthesis~37), 2026. (This archive.)

\bibitem{series:challengecoverage}
Anonymous.
\newblock Challenge Coverage Certificates for Successor Maintenance: Audience Coverage and Support-Only Hooks.
\newblock Series synthesis note (Synthesis~39), 2026. (This archive.)

\bibitem{series:challengeterminalwitnesses}
Anonymous.
\newblock Challenge Terminal Witness Ledgers for Successor Maintenance: Narrow Stop Ownership for Piece Hooks and Support-Only Escalation.
\newblock Series synthesis note (Synthesis~41), 2026. (This archive.)

\bibitem{series:challengeanswerreviewmenus}
Anonymous.
\newblock Challenge Answer Review Menus for Successor Maintenance: Stable Request Classes and Terminal Citation Normal Form.
\newblock Series synthesis note (Synthesis~54), 2026. (This archive.)

\bibitem{series:pairedterminalcitationdiscipline}
Anonymous.
\newblock Paired Terminal Citation Discipline for Review Spines: Stable Joint Reuse of the Checked-Answer Stop and the Request-Side Terminal Stop.
\newblock Series synthesis note (Synthesis~74), 2026. (This archive.)

\bibitem{series:pairedterminalcutoverrules}
Anonymous.
\newblock Paired Terminal Cutover Rules for Review Spines: When Later Terse Papers Stop, Widen, or Reopen.
\newblock Series synthesis note (Synthesis~75), 2026. (This archive.)

\bibitem{series:workedexample}
Anonymous.
\newblock Worked Example for Anonymous-DHT Receipts: Minimal Public Surface, Cross-Release Diff, and Successor Interlock.
\newblock Series synthesis note (Synthesis~17), 2026. (This archive.)

\end{thebibliography}
\end{document}
